\section{Assurance-Contract Failure Modes}
\label{sec:lifecycle-attacks}

This section separates three failure modes that are often conflated.

\textbf{Nominal-state certification.}
A nominal certificate checks safety only at an estimate $\hat x(h)$ or at the reported sensor value. Under FDI, this can accept an action that is safe for $\hat x(h)$ but unsafe for another state in $X_{\mathcal A}(h)$. This is the classical state-semantics failure already addressed by severe-attack safety filters.

\textbf{Action-only filtering.}
A robust safety filter can repair the current action by projecting it into $K_\Gamma(h)$. This protects the physical action applied at that step when the kernel is nonempty. It does not, by itself, certify the learner update. If the update rule still changes $\theta_t$ to $\theta_{t+1}$ such that the raw policy $\pi_{\theta_{t+1}}$ leaves $K_\Gamma(h)$, then the certificate applies to the filtered action, not to the updated learner. The next certificate must therefore reason about the updated action map, not only about the one action that was just projected.

\textbf{Uncertified learning after rejection.}
When $K_\Gamma(h)=\emptyset$, a severe-attack filter can reject or enter fail-safe mode. If the learner nevertheless updates $\theta_t$ using the same attacked history, this update is not certified learning. The mechanism may still operate as an engineering fail-safe, but it cannot claim that the online learner remained inside a sound certificate lifecycle.

\textbf{Protected adaptation, unsafe deployment.}
An adaptation shield can keep every data-collection transition safe while an
attacker corrupts only the log consumed by the learner.  If the resulting
snapshot is later deployed without the shield, the attack succeeds only
through the legitimate update rule.  This delayed failure is distinct from
evading the runtime shield.  A sound snapshot lifecycle therefore requires a
commit certificate before the protected adaptation environment releases a raw
policy.

These attacks motivate separate metrics. An \emph{unsafe certified action}
counts accepted certificates whose applied action is unsafe for some plausible
state. A \emph{stale certified policy} counts accepted certificates for which
the updated raw policy leaves the kernel. \emph{Uncertified learning} counts
nontrivial updates after certificate rejection.  End-to-end studies
add delayed deployment violations, admission coverage, paired clean/poison
outcomes, update fraction, intervention, reward, and certificate runtime.
